\section{Codificaciones Posicionales de Todo el Modelo}
\label{sec:annex-codificaciones-posicionales}

Este anexo documenta el subsistema unificado de codificación posicional
(PE) que aplica una única elección de \texttt{model.positional\_encoding}
a todos los mezcladores de atención de la red. El contrato del esquema es:
\[
\texttt{positional\_encoding} \in \{\texttt{rope},\ \texttt{hope},\ \texttt{nope},\ \texttt{alibi},\ \texttt{bam},\ \texttt{pape},\ \texttt{pape\_efficient},\ \texttt{pape\_ri},\ \texttt{sinusoidal\_absolute},\ \texttt{sinusoidal\_rotary},\ \texttt{learned\_absolute},\ \texttt{none}\}
\]
La implementación reside en \texttt{src/model/embeddings/} y los
ayudantes de inyección compartidos en
\texttt{src/model/attention/common.py}. Las claves bibliográficas citadas
aquí se resuelven contra \texttt{docs/bibliography/}.

\subsection{Arquitectura: Un Módulo Compartido, Exclusión Por Mezclador}

Un único módulo de PE se construye una vez en el codificador
(\texttt{FrankensteinEncoder} o \texttt{FrankensteinViT}) mediante
\texttt{build\_pos\_encoder(config)} y se inyecta en cada
\texttt{HybridLayer}, que lo reenvía a cada mezclador en
\texttt{\_\_init\_\_(config, pos\_encoder=None)} y
\texttt{forward(..., pos\_encoder=None)}.

Cada mezclador lleva una bandera \texttt{use\_pe} por mezclador
(\texttt{<mezclador>\_attn\_use\_pe}, por defecto \texttt{True} para
mezcladores de atención, \texttt{False} para mezcladores
recurrentes/de decaimiento como \texttt{retnet}, \texttt{mamba},
\texttt{ode}, \texttt{gla\_attn}, \texttt{deltanet\_attn},
\texttt{mtla\_attn}). Cuando \texttt{use\_pe} es \texttt{False}, el
mezclador ignora completamente el módulo inyectado.

La PE se aplica en dos fases mediante los ayudantes compartidos
\texttt{apply\_pe\_to\_qk()} (llamado antes del cálculo de puntuaciones de
atención) y \texttt{apply\_pe\_to\_scores()} (llamado después). El
ayudante despacha según el tipo de PE y se convierte en una operación
nula cuando la bandera \texttt{use\_pe} del mezclador es \texttt{False}.

\subsection{Estilos de Inyección}

Los 12 valores del enum se agrupan en cuatro estilos de inyección:

\begin{enumerate}
\item \textbf{Rotación} (\texttt{rope}, \texttt{hope},
   \texttt{sinusoidal\_rotary}): rota los tensores de consulta y clave
   antes del cálculo de puntuaciones. El producto interno
   $\mathbf{q}_i^\top\mathbf{k}_j$ se convierte en función del
   desplazamiento relativo $i-j$.
\item \textbf{Sesgo de puntuación} (\texttt{alibi}): añade un sesgo
   proporcional a la distancia a la matriz de puntuaciones de atención
   \emph{después} de que se calcule.
\item \textbf{Aumentación de Q/K} (\texttt{pape}, \texttt{pape\_efficient},
   \texttt{pape\_ri}): concatena dimensiones posicionales adicionales a
   los vectores de consulta y clave antes del producto punto, aumentando
   la dimensión efectiva de la cabeza.
\item \textbf{Aditivo a nivel de embedding} (\texttt{sinusoidal\_absolute},
   \texttt{learned\_absolute}): se añade directamente a los embeddings de
   tokens antes de la primera capa; ninguna acción por mezclador se
   realiza en el bloque de atención.
\end{enumerate}

\texttt{nope} y \texttt{none} son ambos operaciones nulas: \texttt{nope}
desactiva explícitamente la codificación posicional (el modelo debe
aprender la posición implícitamente), mientras que \texttt{none} es el
módulo nulo devuelto por la fábrica cuando no se desea ninguna PE.

\subsection{Formulaciones Matemáticas}

\subsubsection{RoPE --- Codificación Posicional Rotatoria}

RoPE \citep{su_rope_2024} rota la consulta y clave de cada cabeza un
ángulo proporcional al índice del token. Para dimensión de cabeza
$d_h$ y posición de token $m$, la matriz de rotación
$\mathbf{R}_m \in \mathbb{R}^{d_h \times d_h}$ actúa sobre sub-espacios
2-D $(x_{2i}, x_{2i+1})$:
\[
\mathbf{R}_m = \operatorname{diag}\!\left(
\begin{bmatrix}\cos m\theta_0 & -\sin m\theta_0 \\ \sin m\theta_0 & \cos m\theta_0\end{bmatrix},
\dots,
\begin{bmatrix}\cos m\theta_{d_h/2-1} & -\sin m\theta_{d_h/2-1} \\ \sin m\theta_{d_h/2-1} & \cos m\theta_{d_h/2-1}\end{bmatrix}
\right),
\]
con frecuencias $\theta_i = \texttt{rope\_base}^{-2i/d_h}$. La puntuación
de atención se convierte en
$\mathbf{q}_m^\top \mathbf{R}_{m-n} \mathbf{k}_n$, función del
desplazamiento relativo $m{-}n$. Un factor \texttt{rope\_scaling} $s$
reescala la posición efectiva a $m/s$ para modos de extrapolación de
longitud.

\subsubsection{HoPE --- Codificación Posicional Rotatoria Hiperbólica}

HoPE \citep{dai_hope_2025} generaliza RoPE a la geometría hiperbólica del
modelo de Lorentz. En lugar de rotaciones euclidianas, HoPE aplica
rotaciones de Lorentz construidas a partir de funciones hiperbólicas,
imponiendo un \emph{decaimiento monótono} de los pesos de atención con
el aumento de la distancia entre tokens --- una propiedad que RoPE carece
porque su patrón de atención oscilatorio puede re-surgir en rangos largos.
RoPE se recupera como caso especial. Las perillas \texttt{hope\_base} y
\texttt{hope\_damping} controlan la frecuencia hiperbólica y la tasa de
decaimiento.

\subsubsection{NoPE --- Sin Codificación Posicional}

NoPE \citep{kazemnejad_nope_2023} es la ausencia explícita de cualquier
inyección posicional. El modelo debe descubrir la posición implícitamente
a partir de la máscara causal (decodificador) o de los datos mismos. El
trabajo citado muestra que NoPE supera a los métodos explícitos de PE
(incluyendo RoPE y ALiBi) en benchmarks de generalización de longitud sin
requerir cómputo adicional, porque un decodificador con atención causal
puede codificar implícitamente la posición relativa a través de sus
patrones de atención.

\subsubsection{ALiBi --- Atención con Sesgos Lineales}

ALiBi \citep{press_alibi_2022} no modifica los embeddings ni los vectores
de consulta/clave. En su lugar añade un sesgo lineal estático específico
por cabeza a la matriz de puntuaciones de atención:
\[
\text{puntuación}_{i,j}^{(h)} = \frac{\mathbf{q}_i^{(h)\top}\mathbf{k}_j^{(h)}}{\sqrt{d_h}} - m_h \cdot (i - j), \qquad m_h = \frac{1}{2^{\lfloor h / \lceil H/8 \rceil \rfloor}},
\]
donde $m_h$ es una pendiente geométrica por cabeza $h$ y el sesgo penaliza
las claves distantes. La perilla \texttt{alibi\_num\_heads} controla el
número de cabezas del calendario de pendientes. ALiBi habilita la
extrapolación de longitud: un modelo entrenado en 1024 tokens generaliza
a 2048 sin reentrenamiento.

\subsubsection{BAM --- Mecanismo de Atención Bayesiana}

BAM \citep{bianchessi_bam_2025} reformula la codificación posicional como
un prior probabilístico sobre posiciones y unifica métodos existentes:
NoPE corresponde a un prior Uniforme, ALiBi a un prior
Uniforme$\times$Laplace. La instanciación práctica, GGD-BAM, añade una
matriz de sesgo de posición relativa con Gaussiana Generalizada
$\mathbf{B}$ a los logits de atención, exactamente como ALiBi pero con
una forma no lineal y aprendible:
\[
B_{i,j}^{(h)} = -\bigl(|j - \alpha_h \, i - \mu_h| + \varepsilon\bigr)^{\beta_h}, \qquad \alpha_h = e^{-\theta_\beta^{(h)}\,\theta_\alpha^{(h)}} \ge 0,
\]
con parámetros aprendibles por cabeza $\theta_\alpha, \theta_\beta$ (init
\texttt{bam\_theta\_init}=0, el prior Uniforme) y localización opcional
$\theta_\mu$ (fijado a 0 por defecto vía \texttt{bam\_learn\_mu=false}). El
piso $\varepsilon$ (\texttt{bam\_eps}, por defecto $10^{-5}$) estabiliza la
exponenciación cuando $\beta < 0$, un régimen relajado que suprime el
contexto local y actúa como ``cabeza de recuperación'' de largo alcance.
Casos especiales: $\beta=1, \mu=0$ recupera ALiBi; $\beta=2$ da un prior
Normal; $0<\beta<1$ produce colas más pesadas que Laplace. El paper
aparea BAM con Scalable Softmax (SSMax), un reescalado transversal de
logits $s\cdot\ln(n)$ (con $s$ aprendible por cabeza) que contrarresta el
``attention fading'' y es componible con cualquier PE vía la bandera
\texttt{use\_ssmax} de nivel superior. BAM SSMax logra recuperación precisa
de información a $500\times$ la longitud de contexto de entrenamiento.

\subsubsection{PaPE --- Codificación Posicional Parabólica}

PaPE \citep{ohrstrom_pape_2026} es una codificación centrada en visión que
aumenta la consulta y clave con dimensiones posicionales extra construidas
a partir de funciones base parabólicas. Dados $P$ parábolas
(\texttt{pape\_num\_parabolas}) y $S$ posiciones
(\texttt{pape\_num\_positions}), la característica posicional para la
coordenada $\mathbf{p} \in \mathbb{R}^D$ es
\[
\phi_k(\mathbf{p}) = \sum_{j=1}^{S} \alpha_{k,j}\, f_j(\mathbf{p}), \qquad
f_j(\mathbf{p}) = -\|\mathbf{p} - \mathbf{c}_j\|^2 + b_j,
\]
una suma de $S$ paraboloides centrados en $\mathbf{c}_j$ con
desplazamiento $b_j$. Las $P \cdot S$ dimensiones extra se concatenan a
la consulta y clave antes del producto punto, por lo que la dimensión
efectiva de la cabeza crece de $d_h$ a $d_h + P \cdot S$. Se proveen tres
variantes:

\begin{itemize}
\item \texttt{pape} --- la implementación de referencia.
\item \texttt{pape\_efficient} --- una reformulación en PyTorch puro que
   evita kernels Triton personalizados preservando la misma matemática.
\item \texttt{pape\_ri} --- la variante invariante a rotación
   (PaPE-RI en el paper), que satisface invariancia rotacional por
   construcción y se controla con \texttt{pape\_rotation\_invariant}.
\end{itemize}

PaPE está diseñado para tokens de visión $n$-D (imágenes, video, cámaras
de eventos, nubes de puntos) y es la PE recomendada para la ruta ViT.

\subsubsection{Sinusoidal Absoluta}

La codificación absoluta original del Transformer
\citep{vaswani_attention_2017} añade un vector sinusoidal fijo al
embedding del token:
\[
\text{PE}_{(m,2i)} = \sin\!\left(\frac{m}{\texttt{sinusoidal\_base}^{2i/d}}\right), \quad
\text{PE}_{(m,2i+1)} = \cos\!\left(\frac{m}{\texttt{sinusoidal\_base}^{2i/d}}\right),
\]
con \texttt{sinusoidal\_max\_len} limitando la tabla precomputada y
\texttt{sinusoidal\_scale} multiplicando el embedding. Esta es una PE
aditiva: se aplica una vez a nivel de embedding y los mezcladores no
aplican ninguna PE por capa.

\subsubsection{Sinusoidal Rotatoria}

Una variante rotatoria de la codificación sinusoidal: las mismas
frecuencias sinusoidales se usan para construir una matriz de rotación
aplicada a la consulta y clave (como en RoPE), pero con las perillas
\texttt{sinusoidal\_base} y \texttt{sinusoidal\_scale}. Esta es una PE
de estilo rotacional aplicada por capa mediante
\texttt{apply\_pe\_to\_qk}.

\subsubsection{Absoluta Aprendida}

Una codificación absoluta aprendible: una matriz de parámetros de
$\texttt{learned\_max\_len} \times d$ se añade a los embeddings de
tokens a nivel de embedding. La desviación estándar de inicialización se
controla con \texttt{learned\_init\_std}. Al igual que la sinusoidal
absoluta, los mezcladores no realizan ninguna acción por capa.

\subsection{Banderas \texttt{use\_pe} Por Mezclador}

Cada mezclador expone un campo de configuración
\texttt{<mezclador>\_attn\_use\_pe} (p.\,ej.\
\texttt{standard\_attn\_use\_pe}, \texttt{titan\_attn\_use\_pe}). Los
valores por defecto son:

\begin{itemize}
\item \textbf{True} (mezcladores de atención): \texttt{standard\_attn},
  \texttt{titan\_attn}, \texttt{sigmoid\_attn},
  \texttt{gated\_softmax\_attn}, \texttt{gqa\_attn}, todos los
  mezcladores latentes (\texttt{mla\_attn}, \texttt{gqla\_attn},
  \texttt{mlra\_attn}, \texttt{tucker\_attn}, \texttt{iha\_attn},
  \texttt{gta\_attn}, \texttt{cca\_attn}, \texttt{ccgqa\_attn},
  \texttt{gma\_attn}), todos los mezcladores dispersos
  (\texttt{longformer\_attn}, \texttt{bigbird\_attn},
  \texttt{sparse\_transformer\_attn}, \texttt{sparsek\_attn},
  \texttt{nsa\_attn}, \texttt{msa\_attn}, \texttt{sparda\_attn},
  \texttt{fasa\_attn}, \texttt{sparge\_attn}), y
  \texttt{engram\_attn}.
\item \textbf{False} (mezcladores recurrentes/de decaimiento):
  \texttt{retnet}, \texttt{retnet\_attn}, \texttt{mamba}, \texttt{ode},
  \texttt{gla\_attn}, \texttt{deltanet\_attn},
  \texttt{gated\_deltanet\_attn}, \texttt{gated\_deltanet2\_attn},
  \texttt{hgrn2\_attn}, \texttt{fox\_attn}, \texttt{kda\_attn},
  \texttt{mtla\_attn}.
\end{itemize}

El valor por defecto \texttt{False} para mezcladores recurrentes/de
decaimiento refleja que estas arquitecturas incorporan su propio sesgo
inductivo temporal/posicional en su recurrencia (p.\,ej.\ la rotación
compleja de RetNet, el barrido selectivo de Mamba). Sobrescribir
\texttt{use\_pe=True} está permitido pero es típicamente redundante.

\subsection{Unificación del Transformer de Visión}

La ruta ViT (\texttt{FrankensteinViT}) usaba previamente su propio campo
\texttt{pos\_embedding\_type} con valores \texttt{learned\_1d},
\texttt{none}. Esto se ha unificado con el enum de todo el modelo:
\texttt{learned\_1d} se migra a \texttt{learned\_absolute}, y
\texttt{none} se mapea a \texttt{none}. El ViT construye el mismo módulo
\texttt{build\_pos\_encoder(config)} y lo inyecta en sus HybridLayers,
por lo que un ViT ahora puede usar RoPE, HoPE, ALiBi, PaPE, o cualquier
otra PE del enum.

\subsection{Compatibilidad con Legado}

La bandera booleana legada \texttt{use\_hope} está deprecada y mapeada al
nuevo enum: \texttt{use\_hope=True} $\to$
\texttt{positional\_encoding=hope}, \texttt{use\_hope=False} $\to$
\texttt{positional\_encoding=rope}. La sobrescritura
\texttt{positional\_encoding} por mezclador que previamente existía solo
en \texttt{titan\_attn} es ahora una operación nula: el enum de todo el
modelo es la única fuente de verdad, y la lista de permitidos de
TitanAttention que restringía PE a \texttt{\{hope, rope, none, nope\}} se
ha eliminado --- los 12 valores ahora funcionan con cada mezclador.

\subsection{Perfil Arquitectónico: Codificaciones Posicionales}

La Tabla~\ref{tab:all-pos-enc} resume las 12 codificaciones posicionales
soportadas.

\footnotesize
\begin{longtable}{p{2.4cm}p{1.7cm}p{4.2cm}p{2.0cm}p{2.0cm}}
\toprule
\textbf{Nombre} & \textbf{Estilo} & \textbf{Idea Clave} & \textbf{Parámetros} & \textbf{Referencia} \\
\midrule
\endfirsthead
\toprule
\textbf{Nombre} & \textbf{Estilo} & \textbf{Idea Clave} & \textbf{Parámetros} & \textbf{Referencia} \\
\midrule
\endhead
\midrule
\multicolumn{5}{r}{\textit{continúa en la siguiente página}} \\
\endfoot
\bottomrule
\caption{Catálogo de codificaciones posicionales. ``Estilo'' es el mecanismo de inyección (rotación, sesgo de puntuación, aumentación de Q/K, aditivo). ``Parámetros'' lista las perillas principales de configuración.}
\label{tab:all-pos-enc}
\endlastfoot
rope & Rotación & Rotación basada en frecuencia de Q/K; posición relativa mediante geometría de producto interno & \texttt{rope\_base}, \texttt{rope\_scaling} & \citep{su_rope_2024} \\
hope & Rotación & Rotación de Lorentz hiperbólica; decaimiento monótono de atención con distancia & \texttt{hope\_base}, \texttt{hope\_damping} & \citep{dai_hope_2025} \\
nope & Ninguno & Sin PE explícita; el modelo aprende posición implícitamente de la máscara causal & --- & \citep{kazemnejad_nope_2023} \\
alibi & Sesgo de puntuación & Penalización lineal estática de distancia añadida a puntuaciones de atención; extrapolación de longitud & \texttt{alibi\_num\_heads} & \citep{press_alibi_2022} \\
bam & Sesgo de puntuación & Sesgo de posición relativa con Gaussiana Generalizada; forma ($\beta$) y escala ($\alpha$) aprendibles por cabeza; recupera ALiBi en $\beta=1$; reescalado SSMax vía \texttt{use\_ssmax} & \texttt{bam\_theta\_init}, \texttt{bam\_learn\_mu}, \texttt{bam\_eps}, \texttt{use\_ssmax}, \texttt{ssmax\_s\_init} & \citep{bianchessi_bam_2025} \\
pape & Aum.\ Q/K & Funciones base parabólicas concatenadas a Q/K; centrada en visión, $n$-D & \texttt{pape\_num\_parabolas}, \texttt{pape\_num\_posiciones} & \citep{ohrstrom_pape_2026} \\
pape\_efficient & Aum.\ Q/K & PaPE en PyTorch puro; sin kernels Triton personalizados & igual que pape & \citep{ohrstrom_pape_2026} \\
pape\_ri & Aum.\ Q/K & PaPE invariante a rotación (PaPE-RI) & \texttt{pape\_rotation\_invariant} & \citep{ohrstrom_pape_2026} \\
sinusoidal\_absolute & Aditivo & Sinusoidal fijo añadido a embeddings (Transformer original) & \texttt{sinusoidal\_base}, \texttt{sinusoidal\_max\_len}, \texttt{sinusoidal\_scale} & \citep{vaswani_attention_2017} \\
sinusoidal\_rotary & Rotación & Frecuencias sinusoidales como matriz de rotación en Q/K & igual que sinusoidal & \citep{vaswani_attention_2017} \\
learned\_absolute & Aditivo & Matriz de parámetros aprendible añadida a embeddings & \texttt{learned\_max\_len}, \texttt{learned\_init\_std} & \citep{vaswani_attention_2017} \\
none & Ninguno & Módulo nulo; ninguna información posicional inyectada & --- & --- \\
\bottomrule
\end{longtable}
\normalsize